\documentclass[a4paper,12pt,twoside]{memoir}

% Castellano
\usepackage[spanish,es-tabla]{babel}
\selectlanguage{spanish}
\usepackage[utf8]{inputenc}
\usepackage[T1]{fontenc}
\usepackage{lmodern} % Scalable font
\usepackage{microtype}
\usepackage{placeins}

\RequirePackage{booktabs}
\RequirePackage[table]{xcolor}
\RequirePackage{xtab}
\RequirePackage{multirow}

% Links
\usepackage[colorlinks]{hyperref}
\hypersetup{
	allcolors = {red}
}

% Ecuaciones
\usepackage{amsmath}

% Rutas de fichero / paquete
\newcommand{\ruta}[1]{{\sffamily #1}}

% Párrafos
\nonzeroparskip


% Imagenes
\usepackage{graphicx}
\newcommand{\imagen}[2]{
	\begin{figure}[!h]
		\centering
		\includegraphics[width=0.9\textwidth]{#1}
		\caption{#2}\label{fig:#1}
	\end{figure}
	\FloatBarrier
}

\newcommand{\imagenflotante}[2]{
	\begin{figure}%[!h]
		\centering
		\includegraphics[width=0.9\textwidth]{#1}
		\caption{#2}\label{fig:#1}
	\end{figure}
}



% El comando \figura nos permite insertar figuras comodamente, y utilizando
% siempre el mismo formato. Los parametros son:
% 1 -> Porcentaje del ancho de página que ocupará la figura (de 0 a 1)
% 2 --> Fichero de la imagen
% 3 --> Texto a pie de imagen
% 4 --> Etiqueta (label) para referencias
% 5 --> Opciones que queramos pasarle al \includegraphics
% 6 --> Opciones de posicionamiento a pasarle a \begin{figure}
\newcommand{\figuraConPosicion}[6]{%
  \setlength{\anchoFloat}{#1\textwidth}%
  \addtolength{\anchoFloat}{-4\fboxsep}%
  \setlength{\anchoFigura}{\anchoFloat}%
  \begin{figure}[#6]
    \begin{center}%
      \Ovalbox{%
        \begin{minipage}{\anchoFloat}%
          \begin{center}%
            \includegraphics[width=\anchoFigura,#5]{#2}%
            \caption{#3}%
            \label{#4}%
          \end{center}%
        \end{minipage}
      }%
    \end{center}%
  \end{figure}%
}

%
% Comando para incluir imágenes en formato apaisado (sin marco).
\newcommand{\figuraApaisadaSinMarco}[5]{%
  \begin{figure}%
    \begin{center}%
    \includegraphics[angle=90,height=#1\textheight,#5]{#2}%
    \caption{#3}%
    \label{#4}%
    \end{center}%
  \end{figure}%
}
% Para las tablas
\newcommand{\otoprule}{\midrule [\heavyrulewidth]}
%
% Nuevo comando para tablas pequeñas (menos de una página).
\newcommand{\tablaSmall}[5]{%
 \begin{table}
  \begin{center}
   \rowcolors {2}{gray!35}{}
   \begin{tabular}{#2}
    \toprule
    #4
    \otoprule
    #5
    \bottomrule
   \end{tabular}
   \caption{#1}
   \label{tabla:#3}
  \end{center}
 \end{table}
}

%
% Nuevo comando para tablas pequeñas (menos de una página).
\newcommand{\tablaSmallSinColores}[5]{%
 \begin{table}[H]
  \begin{center}
   \begin{tabular}{#2}
    \toprule
    #4
    \otoprule
    #5
    \bottomrule
   \end{tabular}
   \caption{#1}
   \label{tabla:#3}
  \end{center}
 \end{table}
}

\newcommand{\tablaApaisadaSmall}[5]{%
\begin{landscape}
  \begin{table}
   \begin{center}
    \rowcolors {2}{gray!35}{}
    \begin{tabular}{#2}
     \toprule
     #4
     \otoprule
     #5
     \bottomrule
    \end{tabular}
    \caption{#1}
    \label{tabla:#3}
   \end{center}
  \end{table}
\end{landscape}
}

%
% Nuevo comando para tablas grandes con cabecera y filas alternas coloreadas en gris.
\newcommand{\tabla}[6]{%
  \begin{center}
    \tablefirsthead{
      \toprule
      #5
      \otoprule
    }
    \tablehead{
      \multicolumn{#3}{l}{\small\sl continúa desde la página anterior}\\
      \toprule
      #5
      \otoprule
    }
    \tabletail{
      \hline
      \multicolumn{#3}{r}{\small\sl continúa en la página siguiente}\\
    }
    \tablelasttail{
      \hline
    }
    \bottomcaption{#1}
    \rowcolors {2}{gray!35}{}
    \begin{xtabular}{#2}
      #6
      \bottomrule
    \end{xtabular}
    \label{tabla:#4}
  \end{center}
}

%
% Nuevo comando para tablas grandes con cabecera.
\newcommand{\tablaSinColores}[6]{%
  \begin{center}
    \tablefirsthead{
      \toprule
      #5
      \otoprule
    }
    \tablehead{
      \multicolumn{#3}{l}{\small\sl continúa desde la página anterior}\\
      \toprule
      #5
      \otoprule
    }
    \tabletail{
      \hline
      \multicolumn{#3}{r}{\small\sl continúa en la página siguiente}\\
    }
    \tablelasttail{
      \hline
    }
    \bottomcaption{#1}
    \begin{xtabular}{#2}
      #6
      \bottomrule
    \end{xtabular}
    \label{tabla:#4}
  \end{center}
}

%
% Nuevo comando para tablas grandes sin cabecera.
\newcommand{\tablaSinCabecera}[5]{%
  \begin{center}
    \tablefirsthead{
      \toprule
    }
    \tablehead{
      \multicolumn{#3}{l}{\small\sl continúa desde la página anterior}\\
      \hline
    }
    \tabletail{
      \hline
      \multicolumn{#3}{r}{\small\sl continúa en la página siguiente}\\
    }
    \tablelasttail{
      \hline
    }
    \bottomcaption{#1}
  \begin{xtabular}{#2}
    #5
   \bottomrule
  \end{xtabular}
  \label{tabla:#4}
  \end{center}
}



\definecolor{cgoLight}{HTML}{EEEEEE}
\definecolor{cgoExtralight}{HTML}{FFFFFF}

%
% Nuevo comando para tablas grandes sin cabecera.
\newcommand{\tablaSinCabeceraConBandas}[5]{%
  \begin{center}
    \tablefirsthead{
      \toprule
    }
    \tablehead{
      \multicolumn{#3}{l}{\small\sl continúa desde la página anterior}\\
      \hline
    }
    \tabletail{
      \hline
      \multicolumn{#3}{r}{\small\sl continúa en la página siguiente}\\
    }
    \tablelasttail{
      \hline
    }
    \bottomcaption{#1}
    \rowcolors[]{1}{cgoExtralight}{cgoLight}

  \begin{xtabular}{#2}
    #5
   \bottomrule
  \end{xtabular}
  \label{tabla:#4}
  \end{center}
}


















\graphicspath{ {./img/} }

% Capítulos
\chapterstyle{bianchi}
\newcommand{\capitulo}[2]{
	\setcounter{chapter}{#1}
	\setcounter{section}{0}
	\chapter*{#2}
	\addcontentsline{toc}{chapter}{#2}
	\markboth{#2}{#2}
}

% Apéndices
\renewcommand{\appendixname}{Apéndice}
\renewcommand*\cftappendixname{\appendixname}

\newcommand{\apendice}[1]{
	%\renewcommand{\thechapter}{A}
	\chapter{#1}
}

\renewcommand*\cftappendixname{\appendixname\ }

% Formato de portada
\makeatletter
\usepackage{xcolor}
\newcommand{\tutor}[1]{\def\@tutor{#1}}
\newcommand{\course}[1]{\def\@course{#1}}
\definecolor{cpardoBox}{HTML}{E6E6FF}
\def\maketitle{
  \null
  \thispagestyle{empty}
  % Cabecera ----------------
\noindent\includegraphics[width=\textwidth]{cabecera}\vspace{1cm}%
  \vfill
  % Título proyecto y escudo informática ----------------
  \colorbox{cpardoBox}{%
    \begin{minipage}{.8\textwidth}
      \vspace{.5cm}\Large
      \begin{center}
      \textbf{TFG del Grado en Ingeniería Informática}\vspace{.6cm}\\
      \textbf{\LARGE\@title{}}
      \end{center}
      \vspace{.2cm}
    \end{minipage}

  }%
  \hfill\begin{minipage}{.20\textwidth}
    \includegraphics[width=\textwidth]{escudoInfor}
  \end{minipage}
  \vfill
  % Datos de alumno, curso y tutores ------------------
  \begin{center}%
  {%
    \noindent\LARGE
    Presentado por \@author{}\\ 
    en Universidad de Burgos --- \@date{}\\
    Tutores: \@tutor{}\\
  }%
  \end{center}%
  \null
  \cleardoublepage
  }
\makeatother

\newcommand{\nombre}{Francisco Crespo Diez} %%% cambio de comando

% Datos de portada
\title{Desarrollo de una interfaz para planta piloto}
\author{\nombre}
\tutor{Dr. Alejandro Merino Gómez - Dr. Daniel Sarabia Ortiz}
\date{\today}

\begin{document}

\maketitle


\newpage\null\thispagestyle{empty}\newpage


%%%%%%%%%%%%%%%%%%%%%%%%%%%%%%%%%%%%%%%%%%%%%%%%%%%%%%%%%%%%%%%%%%%%%%%%%%%%%%%%%%%%%%%%
\thispagestyle{empty}


\noindent\includegraphics[width=\textwidth]{cabecera}\vspace{1cm}

\noindent Dr. Daniel Sarabia Ortiz, profesro del departamento de Ingeniería Electromecánica, área de Ingeniería de Sistemas y Automática, y Dr. Alejandro Merino Gómez, profesor del departamento de Ingeniería Electromecánica, área de Ingeniería de Sistemas y Automática.

\noindent Expone:

\noindent Que el alumno D. \nombre, con DNI 71296830G, ha realizado el Trabajo final de Grado en Ingeniería Informática titulado "Desarrollo de una interfaz para planta piloto". 

\noindent Y que dicho trabajo ha sido realizado por el alumno bajo la dirección del que suscribe, en virtud de lo cual se autoriza su presentación y defensa.

\begin{center} %\large
En Burgos, {\large \today}
\end{center}

\vfill\vfill\vfill

% Author and supervisor
\begin{minipage}{0.45\textwidth}
\begin{flushleft} %\large
Vº. Bº. del Tutor:\\[2cm]
D. Daniel Sarabia Ortiz
\end{flushleft}
\end{minipage}
\hfill
\begin{minipage}{0.45\textwidth}
\begin{flushleft} %\large
Vº. Bº. del co-tutor:\\[2cm]
D. Alejandro Merino Gómez
\end{flushleft}
\end{minipage}
\hfill

\vfill

% para casos con solo un tutor comentar lo anterior
% y descomentar lo siguiente
%Vº. Bº. del Tutor:\\[2cm]
%D. nombre tutor


\newpage\null\thispagestyle{empty}\newpage




\frontmatter

% Abstract en castellano
\renewcommand*\abstractname{Resumen}
\begin{abstract}
El objetivo de este proyecto es el desarrollo de una interfaz que actúe como intermediario entre el usuario y una placa Freescale FRDMK64F, con la finalidad de monitorizar variables del proceso que esté contenido en la placa mencionada, así como para actuar sobre dichas variables. 

La placa Freescale FRDMK64F estará a su vez conecatada a una planta piloto en la cual se podrá tener acceso al control de la temperatura y del caudal de aire que ocurre en su interior.
\end{abstract}

\renewcommand*\abstractname{Descriptores}
\begin{abstract}
Interfaz hombre-máquina, monitorización, NXP, Freescale, FRDMK64F, .NET, Windows Forms, C\#\ldots
\end{abstract}

\clearpage

% Abstract en inglés
\renewcommand*\abstractname{Abstract}
\begin{abstract}
The goal of this project is the development of an interface that acts as an intermediary between the user and a Freescale FRDMK64F board, in order to monitor process variables contained in the board mentioned, as well as to act on these variables. 

The Freescale FRDMK64F board will, in turn, be connected to a pilot plant in which it will be possible to have access to the control of the temperature and the flow of air that occurs inside.
\end{abstract}

\renewcommand*\abstractname{Keywords}
\begin{abstract}
Man machine interface, monitoring, NXP, Freescale, FRDMK64F, .NET, Windows Forms, C\#\ldots
\end{abstract}

\clearpage

% Indices
\tableofcontents

\clearpage

\listoffigures

\clearpage

\listoftables
\clearpage

\mainmatter
\capitulo{1}{Introducción}

Desde el inicio del estudio de ingenierías, siempre ha resultado más sencilla la comprensión por parte de los alumnos de conceptos teóricos si estos son vistos en la práctica. Sin embargo, no siempre es posible tener acceso a estas facilidades debido a la no inclusión de un software dirigido al estudiante, entre otras desventajas.

Hasta el momento los alumnos de la Universidad de Burgos han tenido que trabajar con las placas NXP FRDM-K64F de una manera arcaica, con un software de terceros, a través de una consola de comandos que enviaba mensajes directamente a la placa, sin poder optar a una visualización apropiada de los datos.

Partiendo de este principio y teniendo en cuenta los últimas tecnologías implantados en la Universidad de Burgos, el objetivo de este proyecto es el desarrollo de una interfaz hombre-máquina capaz de comunicarse con una placa NXP FRDM-K64F en un entorno amigable, permitiendo al alumno una mayor libertad para trabajar con dicho dispositivo.

\subsection{1.1. Estructura de la memoria}

La memoria sigue la siguiente estructura:
\begin{itemize}
	\item \textbf{Introducción:} breve descripción del problema a resolver y la solución propuesta. Estructura de la memoria y listado de materiales adjuntos.
	\item \textbf{Objetivos del proyecto:} exposición de los objetivos que persigue el proyecto.
	\item \textbf{Conceptos teóricos:} en este apartado se describen los conceptos que se necesitan saber antes de poder comprender todo el proyecto.
	\item \textbf{Técnicas y herramientas:} donde se indican todas las técnicas, bibliotecas, lenguajes, y otras herramientas que se han utilizado durante el proyecto.
	\item \textbf{Aspectos relevantes del desarrollo:} descripción del desarrollo general que ha llevado el proyecto.
	\item \textbf{Trabajos relacionados:} se nombran otros trabajos similares o relacionados con este proyecto. 
	\item \textbf{Conclusiones y líneas de trabajo futuras:} conclusiones obtenidas tras la realización del proyecto y posibilidades de mejora o expansión de la solución aportada.
\end{itemize}

Junto con la memoria se proporcionan los siguientes anexos:
\begin{itemize}
	\item \textbf{Plan del proyecto software:} planificación temporal y estudio de viabilidad del proyecto.
	\item \textbf{Especificación de requisitos del software:} se describe la fase de análisis; los objetivos generales, el catálogo de requisitos del sistema y la especificación de requisitos funcionales y no funcionales.
	\item \textbf{Especificación de diseño:} se describe la fase de diseño; el ámbito
del software, el diseño de datos, el diseño procedimental y el diseño arquitectónico.
	\item \textbf{Manual del programador:} recoge los aspectos más relevantes relacionados
con el código fuente (estructura, compilación, instalación, ejecución, pruebas, etc.).
	\item \textbf{Manual de usuario:} guía de usuario para el correcto manejo de la
aplicación.
\end{itemize}


\subsection{1.2. Materiales adjuntos}
\begin{itemize}
	\item \textbf{Introducción:} 
	\item \textbf{Introducción:} 
	\item \textbf{Introducción:} 
	\item \textbf{Introducción:} 
\end{itemize}
\capitulo{2}{Objetivos del proyecto}

A continuación se detallan los diversos objetivos que han motivado la realización del proyecto.

\section{Objetivos generales}

\begin{itemize}
	\item Crear una interfaz hombre-máquina amigable que permita la monitorización y modificación de los datos gestionados por una placa NXP.
	\item Ayudar a los nuevos estudiantes de la Universidad de Burgos en su comprensión del funcionamiento de una una planta piloto y los sistemas de control necesarios para controlarla. Fundamentalmente orientado para las prácticas de las asignaturas de control en el grado de Electrónica Industrial y Automática, el grado de ingeniería Informática y el máster de Ingeniería Industrial.
	\item Aportar información extra a los datos recibidos por la placa que ayude a la comprensión de los cambios experimentados en la planta piloto.
	\item Dar la posibilidad de almacenar los datos recibidos de la planta piloto de una manera concisa, estructurada y de fácil acceso.
\end{itemize}

\section{Objetivos técnicos}

\begin{itemize}
	\item Desarrollar una aplicación en .NET Framework - Windows Forms para entornos Windows.
	\item Utilizar Git como herramienta de control de versiones distribuido junto con GitHub.
	\item Aplicar la metodología ágil Scrum en el desarrollo del software.
	\item Utilizar la herramienta ZenHub como herramienta de gestión de proyectos.
	\item Documentar a través de la herramienta \LaTeX.
	\item Implementar la comunicación serie entre la aplicación y la la placa para conseguir una comunicación online con la misma y el proceso
	\item Distribuir la aplicación resultante para entornos Windows.
\end{itemize}

\section{Objetivos personales}

\begin{itemize}
	\item Abarcar el máximo número de conocimientos vistos en el grado, como por ejemplo programación orientada a objetos, bases de datos, ingeniería del software, comunicaciones, etc.
	\item Realizar una aportación a la modernización de los recursos de la Universidad de Burgos.
	\item Profundizar en el desarrollo de aplicaciones .NET y en la utilización del entorno de desarrollo Visual Studio.
	\item Trabajar con placas Kinetis de NXP.
	\item Explorar herramientas y metodologías de vanguardia en el mercado laboral.
\end{itemize}
\capitulo{3}{Conceptos teóricos}

En aquellos proyectos que necesiten para su comprensión y desarrollo de unos conceptos teóricos de una determinada materia o de un determinado dominio de conocimiento, debe existir un apartado que sintetice dichos conceptos.

Algunos conceptos teóricos de \LaTeX \footnote{Créditos a los proyectos de Álvaro López Cantero: Configurador de Presupuestos y Roberto Izquierdo Amo: PLQuiz}.

\section{Secciones}

Las secciones se incluyen con el comando section.

\subsection{Subsecciones}

Además de secciones tenemos subsecciones.

\subsubsection{Subsubsecciones}

Y subsecciones. 


\section{Referencias}

Las referencias se incluyen en el texto usando cite \cite{wiki:latex}. Para citar webs, artículos o libros \cite{koza92}.


\section{Imágenes}

Se pueden incluir imágenes con los comandos standard de \LaTeX, pero esta plantilla dispone de comandos propios como por ejemplo el siguiente:

\imagen{escudoInfor}{Autómata para una expresión vacía}



\section{Listas de items}

Existen tres posibilidades:

\begin{itemize}
	\item primer item.
	\item segundo item.
\end{itemize}

\begin{enumerate}
	\item primer item.
	\item segundo item.
\end{enumerate}

\begin{description}
	\item[Primer item] más información sobre el primer item.
	\item[Segundo item] más información sobre el segundo item.
\end{description}
	
\begin{itemize}
\item 
\end{itemize}

\section{Tablas}

Igualmente se pueden usar los comandos específicos de \LaTeX o bien usar alguno de los comandos de la plantilla.

\tablaSmall{Herramientas y tecnologías utilizadas en cada parte del proyecto}{l c c c c}{herramientasportipodeuso}
{ \multicolumn{1}{l}{Herramientas} & App AngularJS & API REST & BD & Memoria \\}{ 
HTML5 & X & & &\\
CSS3 & X & & &\\
BOOTSTRAP & X & & &\\
JavaScript & X & & &\\
AngularJS & X & & &\\
Bower & X & & &\\
PHP & & X & &\\
Karma + Jasmine & X & & &\\
Slim framework & & X & &\\
Idiorm & & X & &\\
Composer & & X & &\\
JSON & X & X & &\\
PhpStorm & X & X & &\\
MySQL & & & X &\\
PhpMyAdmin & & & X &\\
Git + BitBucket & X & X & X & X\\
Mik\TeX{} & & & & X\\
\TeX{}Maker & & & & X\\
Astah & & & & X\\
Balsamiq Mockups & X & & &\\
VersionOne & X & X & X & X\\
} 

\capitulo{4}{Técnicas y herramientas}

En este apartado se detallarán las técnicas y herramientas utilizadas para desarrollar el proyecto.

\section{Metodología}

\subsection{4.1.1. Scrum}

\href{https://proyectosagiles.org/que-es-scrum/}{Scrum} es un marco de trabajo para el desarrollo de \textit{software} que se engloba dentro de las metodologías ágiles. Sus principales características son la adopción de una estrategia de desarrollo incremental, basar la calidad del resultado en el conocimiento tácito de los integrantes de equipos auto organizados, y en el solapamiento de las diferentes fases del desarrollo \cite{web:scrum}.

En este caso los \textit{sprints} han sido de dos semanas, tras los cuales se realizaba una reunión para definir las nuevas tareas del siguiente \textit{sprint} y realizar una retroalimentación del \textit{sprint} recién hecho.

\subsection{4.1.2. Técnica de Pomodoro}

Para mejorar la concentración durante el desarrollo del proyecto se ha utilizado la técnica Pomodoro, método de gestión del tiempo consistente en intervalos de 25 minutos de concentración intensa seguidos de descansos de 5 minutos, y descansos de 30 minutos después de haber completado cuatro ciclos 25+5 \cite{web:pomodoro}.

\section{Patrones de diseño}

\subsection{4.2.1. MVC - Modelo Vista Controlador}

A la hora de estructurar el desarrollo de la aplicación se ha optado por la utilización del Modelo Vista Controlador, consiguiendo así separar la capa que representa la realidad, la capa que conoce los métodos y atributos del modelo, recibiendo y realizando las peticiones del usuario, y la capa visible para el usuario \cite{web:patronDis}.

\section{Control de versiones}

\subsection{4.3.1. Git}

\href{https://git-scm.com/}{Git} es, a día de hoy, el sistema de control de versiones distribuido más usado del mundo. Su propósito es llevar el registro de los cambios realizados en los archivos de un ordenador y coordinar el trabajo de un equipo formado por varias personas sobre un mismo archivo \cite{web:Git}.

La otra opción tenida en cuenta para este proyecto fue \href{https://subversion.apache.org/}{Subversion}, que fue descartada por la facilidad de uso que ofrece Git, la familiaridad con el sistema de control de versiones y las herramientas con las que se puede complementar.

\section{Hosting de repositorio}

\subsection{4.4.1. GitHub}

Una de las plataformas de desarrollo colaborativo más popular cuyo objetivo es alojar proyectos usando el sistema de control de versiones Git es \href{https://github.com/}{GitHub}. Ofrece un repositorio Git donde los desarrolladores pueden almacenar, compartir, testear y colaborar en proyectos web.

Al comenzar el proyecto también se tuvo en cuenta la posibilidad de utilizar \href{https://about.gitlab.com/}{GitLab} pero fue descartado por la compatibilidad de GitHub con otras herramientas usadas en el desarrollo del proyecto y por la familiaridad con dicho hosting de repositorio.

\section{Gestión del proyecto }

\subsection{4.5.1. ZenHub}

\href{https://www.zenhub.com/}{ZenHub} es una herramienta gratuita para proyectos pequeños u \textit{open source} cuya finalidad es la gestión de proyectos en GitHub. Esta herramienta trabaja de manera totalmente integrada ofreciendo un tablero canvas en el que cada \textit{issue} nativo de GitHub se representa como una tarea. A estas tareas se le pueden asignar prioridad, dependencias, estimaciones de tiempo de realización, colaboradores y el \textit{sprint} al que pertenece. ZenHub también permite graficar los avances llevados a cabo en el proyecto.

\imagen{zenhubOverview}{ZenHub Overview}

En un principio también se barajó utilizar GitHub Proyect pero tras ver los vídeos explicativos subidos en la plataforma de la universidad me decanté por ZenHub.

\section{Entorno de desarrollo integrado (IDE)}

\subsection{4.6.1. Microsoft Visual Studio 2017}

\href{https://visualstudio.microsoft.com/es/vs/}{Microsoft Visual Studio 2017} es un entorno de desarrollo integrado para sistemas operativos Windows y macOs con todas las características para Android, iOS, Windows, la Wev y la nube. Este IDE permite escribir código de una manera eficiente y precisa sin perder el contexto del archivo, dotando al programador de facilidades en la programación \cite{web:visualstudio}.

\imagen{visualStudio}{Ventana de desarrolo del proyecto en Visual Studio 2017}

Cuando se inició el proyecto también se tuvo en cuenta usar versiones anteriores de Visual Studio por sus compatibilidades a la hora de desarrollar Windows Forms sobre C++, pero tras enfrentar las ventajas de cada IDE se decidió usar el Visual Studio 2017 con Windows Forms en C\#.

\section{Comunicación}

La comunicación con los tutores se ha realizado a través de varias vías:
\begin{itemize}
	\item Correo electrónico.
	\item Reuniones presenciales. 
	\item Canvas de ZenHub, donde se puede comentar cada tarea.
\end{itemize}

También se pensó en utilizar herramientas como \href{https://slack.com/intl/es-es/}{Slack} pero se descartaron porque se vieron innecesarias.

\section{Documentación}

\subsection{4.8.1. LaTeX}

\LaTeX es un sistema de composición de textos de alta calidad distribuido como software libre, el cual incluye características diseñadas para la producción tanto técnica como científica, siendo así el estándar de facto para la publicación y comunicación de documentos de carácter científico. \cite{web:latex}.

Para la correcta presentación de esta documentación que está siendo leída se ha utilizado la plantilla ofrecida por la Universidad de Burgos, con la finalidad de alcanzar una mayor afinidad en el formato de la misma.

En un principio se planteó la posibilidad de utilizar tanto Open Office Writter como Microsoft Word, en el primer caso porque la Universidad de Burgos ofrecía una plantilla oficial, y en el segundo porque la universidad incluye una licencia para su uso, sin olvidar lo familiarizado que estoy a ella.

\subsection{4.8.2. HTML Help Workshop}

\href{https://www.microsoft.com/en-us/download/details.aspx?id=21138}{HTML Help Workshop} es un programa que permite crear, editar y compilar archivos HTML para convertirlos en proyectos CHM. Este archivo CHM será el encargado de mostrar la ayuda presente en todas las ventanas de la aplicación.

\imagen{chmFile}{CHM File}

\section{Test}

\subsection{4.9.1. Codacy}

Para llevar a cabo una revisión del código de la aplicación se utilizó \href{https://www.codacy.com/}{Codacy}, herramienta que ejecuta de manera automática exámenes sobre el código de cada \textit{commit} que se realice en Git. Codacy informa de los problemas de seguridad, de cobertura, de duplicidad y de complicidad que pueda tener el código, catalogando a este según los baremos de la propia aplicación y graficando los resultados.

\imagen{codacyDashboard}{Codacy Dashboard}

\section{Otras herramientas}

\subsection{4.10.1. Paint.NET}

\href{https://www.getpaint.net/index.html}{Paint.NET} es una herramienta de edición de imágenes sobre entorno de Windows. Ofrece una gran cantidad de posibilidades a la hora de editar una imagen, es un programa ligero y de muy fácil utilización. También se tuvo en cuenta \href{http://www.gimp.org.es/}{Gimp} como herramienta de edición de imágenes pero se descartó porque tardaba mucho en ejecutarse.

\subsection{4.10.2. Termite}

\href{https://www.compuphase.com/software_termite.htm}{Termite} es un terminal \href{https://es.wikipedia.org/wiki/RS-232}{RS232} que actúa como una interfaz de chat teniendo un como canal un puerto serie. Cabe destacar la simplicidad de esta herramienta y su facilidad de instalación. Su uso permitió la comprobación de la comunicación con la placa y la realización de pruebas sobre la misma.

\subsection{4.10.3. Active Presenter}

\href{https://atomisystems.com/activepresenter/}{Active Presenter} es un programa dedicado a la grabación de pantalla, edición de vídeo y creación de tutoriales. Antes de elegir este producto se tuvieron presentes otras opciones como \href{https://www.techsmith.com/video-editor.html}{Camtasia} o \href{https://icecreamapps.com/Screen-Recorder/}{Ice Cream Screen Recorder} pero se eligió Active Presenter por haber trabajado con ello previamente y por las herramientas de edición que ofrece.

\section{Base de datos}

\subsection{4.11.1. Microsoft SQL Server}

\href{https://www.microsoft.com/es-es/sql-server/sql-server-2017}{Microsoft SQL Server} es un sistema de gestión de bases de datos relacionales, diseñado por Microsoft, que se ejecuta en T-SQL (Transact SQL), consiguiendo así añadir características varias a SQL estándar (control de transacciones, excepción y manejo de errores, procesamiento fila, entre otras) \cite{web:sqlServer}.

Dentro de las posibilidades que ofrecía el mercado para cubrir la necesidad de base de datos que tenía la aplicación se tuvieron en cuenta \href{https://mariadb.org/}{MariaDB} y \href{https://www.mysql.com/}{MySQL}. La decisión de usar Microsoft SQL Server se vio motivada a la compatibilidad con el resto de aplicaciones que se estaban usando y por la facilidad de uso que presentaba.

\subsection{4.11.4. Microsoft SQL Server Management Studio 17}

\href{https://docs.microsoft.com/es-es/sql/ssms/download-sql-server-management-studio-ssms?view=sql-server-2017}{Microsoft SQL Server Management Studio 17} es un entorno integrado que permite administrar cualquier infraestructura de SQL.

Se tuvo presente esta herramienta durante el desarrollo de la aplicación para poder interactuar con la base de datos.

\section{Editores}

\subsection{4.12.1. Texmaker}

\href{http://www.xm1math.net/texmaker/}{Texmaker} es un editor multiplataforma gratuito (licencia GNU GLP v3.0) para \LaTeX que integra la mayoría de herramientas necesarias para la escritura de documentos, corrector ortográfico, auto-completado, resultado de análisis, visor de PDF integrado, entre otros.

En un principio se trató de utilizar \href{https://code.visualstudio.com/}{Visual Studio Code} con la extensión LaTeX Workshop pero no ofreció un uso eficiente a la hora de compilar toda la solución \LaTeX.

\subsection{4.12.2. Visual Studio Code}

\href{https://code.visualstudio.com/}{Visual Studio Code} es un editor de código fuente desarrollado por Microsoft para Windows, Linux y MacOs. Visual Studio Code se basa en \href{https://electronjs.org/}{Electron}, un framework que se utiliza para implementar aplicaciones \href{https://nodejs.org/es/}{Node.js} para el escritorio. Este editor permite trabajar no sólo con C\# y VB (lenguajes propietarios de Microsoft), sino con toda una colección de lenguajes Open Source (Java, Python, PHP, JavaScript, JSON, entre otros) \cite{web:vsCode}.

Aparte de incluir IntelliSense (predicciones ofrecidas por el editor mientras el usuario está escribiendo) y control de versiones a través de Git, el software permite instalar un serie de plugins que completan las necesidades que puedan aparecer según el lenguaje con el que se esté programando.

Este editor fue utilizado para el diseño de las páginas HTML usadas en los archivos de ayuda CHM.
\capitulo{5}{Aspectos relevantes del desarrollo del proyecto}

Este apartado pretende recoger los aspectos más importantes del desarrollo del proyecto, presentando los problemas que fueron aconteciendo y los caminos que se eligieron para solventarlos.

\section{Inicio del proyecto}

La elección de este proyecto se vio motivada tanto por el interés personal por desarrollar una aplicación que se comunicara con un dispositivo como por la posibilidad de ayudar a la Universidad de Burgos y a sus alumnos a tener una facilidad más con la que pudieran contar a la hora de llevar a cabo sus estudios y que no fuera un proyecto que queda olvidado.

Durante las primeras tomas de contacto con los tutores, se barajaron las posibilidades en las que el proyecto podría ser desarrollado, eligiendo en un principio trabajar en C++. Esta decisión se vio motivada porque la placa viene programada en C++ y por los conocimientos de los tutores en este lenguaje. Tras un pequeño periodo de formación en C++ y tras investigar las posibilidades que ofrecía C\# para cubrir las necesidades del proyecto, se replanteó el desarrollo del mismo, decidiendo cambiar a C\# y a la última versión de Visual Studio en la que me encontraba más cómodo.

Una vez definido el lenguaje de desarrollo se trataron los temas relacionados con el repositorio que se iba a utilizar, en este caso GitHub, las herramientas y entornos de desarrollo que se iban a utilizar (ZenHub, Visual Studio 2017), y la periocidad con la que se iban a realizar reuniones para tratar los avances del proyecto, siendo periodos de dos semanas.

\section{Formación}

El proyecto requería de unos conocimientos de los que no se disponía originalmente, por lo que hubo que solventar estas cadencias.
\capitulo{6}{Descripción de la herramienta}

También añadir capítulo "Descripción de la herramiento": descripción, funcionalidad y características...

\section{Descripción}



\section{Funcionalidad}



\section{Características}
\capitulo{7}{Trabajos relacionados}

Este apartado sería parecido a un estado del arte de una tesis o tesina. En un trabajo final grado no parece obligada su presencia, aunque se puede dejar a juicio del tutor el incluir un pequeño resumen comentado de los trabajos y proyectos ya realizados en el campo del proyecto en curso. 

\capitulo{8}{Conclusiones y Líneas de trabajo futuras}

En esta sección se exponen tanto las conclusiones derivadas del trabajo, como las posibles líneas de trabajo futuras por las que se puede dar continuidad al proyecto. 

\section{Conclusiones}

Desde un punto de vista general y de una manera personal, finalizo este trabajo de fin de grado llevándome conmigo una sensación positiva, tanto por todo lo que he aprendido gracias a la resolución de los diferentes problemas que me han aparecido, como por dejar una herramienta útil  para la formación de futuros alumnos de la Universidad de Burgos. Poder haber participado en este proyecto comenzado por mis compañeros Rubén Zambrana Rodríguez y María Isabel Revilla Izquierdo me hace sentir que formo parte de una idea mayor que contiene una finalidad altruista y formativa.

Ciñéndome a todo lo relacionado a la documentación, resulta interesante ver como todo aquello aprendido durante la carrera sobre diagramas, previsiones y un largo etcétera de conceptos, es útil para el planteamiento de un proyecto, ayudando al alumno a ponerse en una situación real de desarrollo de software y en la que puede ver todo lo que está involucrado.

Si atiendo a la parte técnica, he tenido la suerte de poder desarrollar la aplicación en el entorno y lenguaje a los que más habituado estoy, por lo que gracias a ello he conseguido ampliar mis conocimientos sobre C\#, Visual Studio y SQL Server, preparándome un poco más para el mundo laboral. Cabe destacar que a la vez que realizaba este trabajo de fin de grado estaba acudiendo a prácticas extracurriculares en la empresa Keyland Sistemas de Gestión, de donde he obtenido muchas ideas y ayuda tanto para organizar el proyecto como para solventar los problemas que han ido aconteciendo.

En resumen, puedo concluir orgulloso que ha sido un proyecto con el que me he nutrido tanto personal como profesionalmente. Durante este tiempo he sido capaz de medirme y ver hasta dónde alcanzan tanto mis conocimientos como mi habilidad para plasmarlos en un entorno real.

\section{Líneas de trabajo futuras}

Durante la programación de la aplicación aparecieron posibles mejoras que se pueden añadir a la solución en líneas de trabajo futuras.

\subsection{Comunicación a través del puerto Ethernet}

Atendiendo a las posibilidades que ofrece la placa Freescale Freedom NXP FRDM-K64F resultaría interesante ampliar las posibilidades de comunicación a través  del puerto Ethernet. Esta línea de trabajo se ha contemplado desde el inicio del proyecto y se puede ver reflejado en la opción de menú de ''Comunicaciones'' de la ventana principal, la cual ofrece la opción de ''Otros''.

\subsection{Internacionalización}

Atendiendo a que el propósito de esta aplicación es educativo y se presenta en una universidad que tiene presencia a nivel internacional, resultaría provechoso añadir tantos idiomas como fueran posibles a la misma, ampliando así las facilidades dadas a los alumnos extranjeros.

Esta línea de trabajo se tuvo en cuenta desde el diseño de la ventana principal, dejando preparado un archivo de recursos con dos idiomas añadidos y enlazando todas las cadenas de texto de la aplicación con dicho archivo.

\subsection{TDD}

En mi opinión creo que resultaría producente realizar las futuras mejoras utilizando TDD (Test-Driven Development), consiguiendo de esta manera una mayor eficiencia y ahorrando tiempo a la larga en la resolución de futuros problemas que pudieran aparecer.

\subsection{Variables PID}

Actualmente la aplicación trabaja exclusivamente con variables normales, y no atiende a que éstas puedan ser de tipo PID, ni incluye la interfaz específica para ellas. La adición de esta mejora completaría el apartado relativo a las variables contenidas en un proyecto y abriría una abanico de funcionalidades nuevas a incluir a su vez.

\subsection{Estilos}

En esta versión la aplicación se muestra con los estilos predeterminados de una aplicación de Windows Forms. Podría resultar visualmente ventajoso trabajar sobre unos estilos con los que dotar a las ventanas, así como efectos y eventos al realizar acciones como el guardado de una variable o la selección de una imagen.



\bibliographystyle{plain}
\bibliography{bibliografia}

\end{document}
